\documentclass[11pt,a4paper]{article}
\usepackage[T1]{fontenc}
\usepackage{lmodern,amsmath,amssymb,amsthm,mathtools,booktabs,microtype}
\usepackage[margin=27mm]{geometry}
\usepackage[colorlinks=true,linkcolor=blue,urlcolor=blue,citecolor=blue]{hyperref}
\newtheorem{theorem}{Theorem}[section]
\newtheorem{lemma}[theorem]{Lemma}
\newtheorem{corollary}[theorem]{Corollary}
\newtheorem{proposition}[theorem]{Proposition}
\newcommand{\Q}{\mathbb Q}\newcommand{\Z}{\mathbb Z}
\DeclareMathOperator{\Res}{Res}\DeclareMathOperator{\Disc}{Disc}
\setlength{\emergencystretch}{2em}
\title{Separating the Nonzero Part of a Finite Algebra}
\author{}\date{9 November 2026}
\begin{document}\maketitle
\begin{abstract}
We give an explicit polynomial, using one additional rational coordinate,
for a guarded rational point of a finite algebra obtained by adjoining
quadratic roots to an algebra with separable fibers. A characteristic
polynomial identifies the factors on which the guard is nonzero.
The product of its nonzero eigenvalues then bounds every required
quadratic root separation; only the scalar discriminants of the initial
separable polynomials are needed. The construction covers degree drops,
nonreal points, and nonreduced fibers where the guard vanishes.
Applied to the specified arithmetic sources, it gives a family in
27 total variables of degree less than $2^{87}$ and a fully specified
polynomial in 29 total variables of degree less than $2^{89}$ and
coefficient height at most $2^{2^{99}}$.
The elimination theorem is unconditional. Assertions about defining the
integers retain the arithmetic premises and the effective or ineffective
thresholds of the respective sources.
\end{abstract}

\section{The algebra and its guarded rational points}

Let $R_0=\Z[x_1,\ldots,x_a]$ and consider the finite free algebras
\[
 B=R_0[z_1,\ldots,z_p]/(f_1(z_1),\ldots,f_p(z_p)),\qquad
 A=B[q_1,\ldots,q_s]/(q_j^2-c_j)_{j=1}^s.
\]
Here $p,s\ge1$, the polynomials $f_i\in R_0[Z]$ are monic
of degrees $d_i\ge2$, and $c_j\in B$. Assume that every $f_i$
is separable at every rational specialization of the base coordinates.
Write $m=p+s$ and $\theta_1,\ldots,\theta_m$ for the generators
in this order, with $d_{p+j}=2$. Put
\[
 b=\prod_{i=1}^p d_i,\qquad n=b2^s.
\]
The monomials with each exponent below its defining degree are
bases of $B$ and $A$. For an element $u$, let $M_u$ denote its
multiplication matrix in the indicated basis.

Fix $g\in B$ and put
\[
 G=gq_1\cdots q_s,\qquad T=g^2c_1\cdots c_s\in B.
\]
Thus $G^2=T$ in $A$. We encode, at a rational base point,
\begin{equation}\label{eq:predicate}
 \text{there is a homomorphism } A\longrightarrow\Q
 \text{ over that base point with }G\ne0.
\end{equation}
There are no further residual equations. Every algebra generator
will be eliminated; precisely one new rational coordinate $U$ is used.
At a rational base point $B$ is an \'etale $\Q$-algebra. On
its factors where $T\ne0$, all $c_j$ and $g$ are nonzero,
so the corresponding part of $A$ is also \'etale. The remaining
part of $A$ may be nonreduced.

\section{An explicit integer polynomial}

Define scalar integer polynomials
\[
 \det_B(I-vM_T)=\sum_{j=0}^b a_jv^j,\quad a_0=1,
 \qquad \delta_i=\Disc(f_i)\quad(1\le i\le p).
\]
Use the convention
$\Disc(f)=(-1)^{d(d-1)/2}\Res(f,f')$ for a monic polynomial
of degree $d$. Define
\begin{equation}\label{eq:rootbound}
 R=n\left(1+\sum_{i=1}^m\sum_{u,v=1}^n
                 (M_{\theta_i}^{A})_{uv}^2
            +\sum_{u,v=1}^b(M_g^{B})_{uv}^2\right),
 \qquad \sigma=\sum_{i=1}^p d_i(d_i-1).
\end{equation}
The superscripts specify the algebra in which multiplication is taken.

For each integer $1\le r\le b$, use only the following polynomial
operations:
\begin{align*}
 c&=a_r,& h_r(Z)&=\sum_{j=0}^r a_jZ^{r-j},& E&=c-h_r(T)\in B,\\
 D&=c^2\prod_{i=1}^p\delta_i^2,& S_r&=r(s+1)+\sigma,\\
 K&=(2R)^{S_r}(1+c^2)\prod_{i=1}^p(1+\delta_i^2),& C&=D+4RK.
\end{align*}
These symbols depend on $r$. Set
\begin{equation}\label{eq:beta}
 \beta_r=DE\left(C^m+\sum_{i=1}^mD^iC^{m-i}\theta_i\right)\in A,
 \qquad H_r(\mathbf x,U)=\det_A(I-UM_{\beta_r}).
\end{equation}
Finally define
\begin{equation}\label{eq:product}
 Q_r=\begin{cases}
 H_r^2+\displaystyle\sum_{j=r+1}^b a_j^2,&r<b,\\
 H_b,&r=b,
 \end{cases}
 \qquad P(\mathbf x,U)=\prod_{r=1}^b Q_r.
\end{equation}
All finite index ranges and integer coefficient operations are explicit.
No division occurs; the expressions remain polynomials where $c=0$.

\begin{theorem}\label{thm:main}
For every rational base point $\mathbf x$, predicate
\eqref{eq:predicate} holds if and only if there is a rational
$U$ satisfying $P(\mathbf x,U)=0$.
\end{theorem}

\section{Proof of the projection theorem}

Fix a rational base point throughout this section. A geometric
character means an algebra homomorphism to $\overline\Q$ over $\Q$.
Let $r_0$ be the number of characters $\chi$ of $B$ with
$T(\chi)\ne0$, counted individually.

\begin{lemma}\label{lem:idempotent}
The degree of $\det_B(I-vM_T)$ is $r_0$.
For $r=r_0>0$, one has $c\ne0$, and $E/c$ equals one
on the factors with $T\ne0$ and zero on all remaining factors.
\end{lemma}
\begin{proof}
Since $B$ is \'etale, the determinant is
$\prod_\chi(1-vT(\chi))$. For $r=r_0$, the polynomial $h_r$
is the product of $Z-T(\chi)$ over those characters with
$T(\chi)\ne0$, including repetitions of their values.
It vanishes at every such value and has constant term $c$.
Thus $E=c-h_r(T)$ has the stated values.
\end{proof}

For the correct case $r=r_0>0$, all scalar $\delta_i$ are
nonzero by hypothesis, so $D>0$, $K>0$, and
\[
 \varepsilon=D/K
 =\frac{c^2}{1+c^2}\prod_{i=1}^p\frac{\delta_i^2}{1+\delta_i^2}
                      (2R)^{-S_r}\in(0,1].
\]

\begin{lemma}\label{lem:separate}
In the correct case, two characters of $A$ with $G\ne0$
which first differ at $\theta_i$ differ there in complex absolute
value by at least $\varepsilon$. The element
\[
 \alpha=L^m+\sum_{i=1}^mL^{m-i}\theta_i,
                 \qquad L=C/D=1+4R/\varepsilon,
\]
has distinct nonzero values on all these characters.
\end{lemma}
\begin{proof}
Every value of every $\theta_i$, and every value of $g$, has
absolute value at most $R$. Indeed these values are eigenvalues
of the respective multiplication matrices. The spectral radius
of a $k$-by-$k$ matrix is at most $k$ times the maximum absolute
entry; $k\le n$ and $|t|\le1+t^2$ prove the assertion from
\eqref{eq:rootbound}. In particular $R\ge1$.

First suppose $i\le p$, and let $\Delta$ be the difference
in the $i$th generator. The discriminant $\delta_i$ is the
product of the squared pairwise root differences of $f_i$.
Every root has absolute value at most $R$. With
$k_i=d_i(d_i-1)$, this gives
\[
 |\Delta|\ge\frac{\sqrt{|\delta_i|}}{(2R)^{(k_i-2)/2}}.
\]
For every real $z>0$ one has $\min(1,z^2)\le\sqrt z$.
Consequently
\[
 \frac{\delta_i^2}{1+\delta_i^2}(2R)^{-k_i}
 \le\frac{\sqrt{|\delta_i|}}{(2R)^{(k_i-2)/2}}.
\]
The other factors in $\varepsilon$ are at most one, and
$S_r\ge k_i$. Hence $\varepsilon\le|\Delta|$.

Next suppose $i=p+j$. Since the earlier point generators agree,
the two characters lie above the same character $\chi$ of $B$,
and their values are the two roots of $Z^2-c_j(\chi)$.
Now
\[
 |c|=\prod_{\psi:\,T(\psi)\ne0}
                  |g(\psi)|^2\prod_{k=1}^s|c_k(\psi)|.
\]
There are $r(s+1)$ factors in this product; every factor is
at most $R^2$. Isolating the particular nonzero factor
$|c_j(\chi)|$ yields
\[
 |c_j(\chi)|\ge |c|R^{-2r(s+1)+2},\qquad
 |\Delta|=2\sqrt{|c_j(\chi)|}
              \ge2\sqrt{|c|}R^{1-r(s+1)}.
\]
On the other hand,
\[
 \varepsilon\le\frac{c^2}{1+c^2}(2R)^{-r(s+1)}
       \le\sqrt{|c|}(2R)^{-r(s+1)}\le|\Delta|.
\]
This proves the separation estimate, also for complex roots.
It uses the product of the nonzero values of $T$, without
requiring those values to be pairwise distinct.

For two characters first differing at $\theta_i$, the leading
term in their $\alpha$-difference has absolute value at least
$\varepsilon L^{m-i}$. The remaining terms have total absolute
value less than
\[
 \frac{2R}{L-1}L^{m-i}=\frac\varepsilon2L^{m-i}.
\]
The values of $\alpha$ are therefore distinct. Also, the sum
after $L^m$ has absolute value less than
$RL^m/(L-1)\le L^m/4$, so none of its values is zero.
\end{proof}

\begin{proof}[Proof of Theorem~\ref{thm:main}]
If $c=0$ in a case, then $D=0$, $\beta_r=0$, and $H_r=1$.
For $r<b$, a rational zero of $Q_r$ therefore forces $c\ne0$
and $a_j=0$ for every $j>r$; it also forces $H_r=0$ by the
sum-of-squares identity. For $r=b$, a zero of $H_b$ similarly
forces $c=a_b\ne0$. Thus any zero uses precisely the correct
positive degree case.

By Lemma~\ref{lem:idempotent}, multiplication by $E$ vanishes
identically on the part of $A$ above $T=0$, including its
nilpotents. On the other part one has
\[
                         \beta_r=cD^{m+1}\alpha.
\]
This part is \'etale. Lemma~\ref{lem:separate} implies that its
$\beta_r$-values are distinct and nonzero. Hence $H_r$ is the
product of the distinct factors $1-U\beta_r(\chi)$ over the
characters with $G\ne0$. A rational zero gives a one-dimensional
rational kernel of $I-UM_{\beta_r}$. All native multiplication
matrices commute with $M_{\beta_r}$ and preserve this line.
Their actions on it are rational scalars satisfying the algebra
relations. They yield a rational character. Its $\beta_r$-value
is nonzero, so it has $G\ne0$. This proves soundness.

Conversely a guarded rational character implies $r_0>0$.
In the correct case its $\beta_r$-value is nonzero and rational.
Its reciprocal annihilates $H_r$, and all tail coefficients
$a_j$, $j>r$, vanish. It therefore annihilates $Q_r$ and $P$.
If $r_0=0$, every $c=a_r$ is zero and every $H_r$ is one,
so no factor has a zero. This includes the empty guarded part.
\end{proof}

\section{Degree and coefficient estimates}

Give each base coordinate weight one, and give the native generators
weights $w_i\ge1$ for which monic reduction does not increase
weight. For coefficient norms use the weighted column norm,
with basis scaling $2^{\sum e_i\gamma_i}$ on the monomial
$\prod\theta_i^{e_i}$. Scalar polynomials have the coefficient
$\ell^1$ norm. Suppose the following bounds are supplied:
\[
 \|M_{\theta_i}\|\le2^{\gamma_i},\quad
 \operatorname{wt}(G)\le w_G,\quad \|M_G\|\le2^\eta,
 \quad \operatorname{wt}(g)\le w_g,\quad\|M_g^B\|\le2^{\eta_g}.
\]
The $\gamma_i,\eta,\eta_g$ are nonnegative integers. The
companion-column recurrence in~\cite{Weighted} supplies these
bounds from the integer coefficient circuit. The $B$-basis uses
the restrictions of the same weights. Since $G^2=T\in B$, its
weight and logarithmic operator norm are bounded by
$t=2w_G$ and $\mu=2\eta$. For the norm assertion, $B$ is an
invariant subspace of multiplication by $G^2$.

For positive integers $u$, write $\ell(u)=\lceil\log_2u\rceil$.
Set
\begin{align*}
 w_*&=\max\left\{\sum_{i=1}^m(d_i-1)w_i+\max_iw_i,
                   \sum_{i=1}^p(d_i-1)w_i+w_g\right\},\\
 \gamma_*&=\max\left\{\sum_{i=1}^m(d_i-1)\gamma_i+\max_i\gamma_i,
                   \sum_{i=1}^p(d_i-1)\gamma_i+\eta_g\right\},\\
 R_d&=2w_*,\qquad
 R_l=\ell(n)+\ell(1+mn^2+b^2)+2\gamma_*,\qquad a_l=\mu+\ell(b),\\
 u_i&=d_i(d_i-1)w_i,\qquad
 v_i=d_i(d_i-1)\gamma_i+
                    \ell\bigl((2d_i-1)!d_i^{d_i}\bigr)\quad(i\le p),\\
 u&=\sum_{i=1}^p u_i,\qquad v=\sum_{i=1}^p v_i.
\end{align*}
Define four integer majorants
\begin{align*}
 W_1&=(3+2m)t+m(s+1)R_d,\\
 W_0&=2(m+1)u+m(\sigma+1)R_d+\max_iw_i,\\
 L_1&=(3+2m)a_l+m(s+1)(1+R_l),\\
 L_0&=2(m+1)v+\ell(b)+\ell(m+1)\\
    &\hspace{8mm}+m\bigl(\sigma(1+R_l)+p+4+R_l\bigr)+\max_i\gamma_i.
\end{align*}
The letters $L_0,L_1$ here denote coefficient bounds and are
unrelated to the fixed coefficient exponent in the applications.

\begin{theorem}\label{thm:bounds}
The polynomial in~\eqref{eq:product} satisfies
\begin{align}
 \deg P&\le n\bigl(W_1b^2+(W_0+1)(2b-1)\bigr),\label{eq:degree}\\
 \log_2\|P\|_1&\le
 n\bigl(L_1b^2+(L_0+\ell(n)+1)(2b-1)\bigr)
                  +(b-1)\ell(b+1).\label{eq:height}
\end{align}
Its coefficient height is at most $\|P\|_1$.
\end{theorem}
\begin{proof}
An entry in the multiplication matrix of an element of weight $w$
has degree at most $w+$ column basis weight $-$ row basis weight.
The analogous logarithmic coefficient-norm estimate uses the
$\gamma$ weights. This proves the bounds $R_d,R_l$ for
\eqref{eq:rootbound}.

In every principal minor the row and column weights cancel.
If multiplication by an element of an algebra of rank $h$ has
weight at most $w$ and operator norm at most $2^\lambda$, then
the $j$th characteristic coefficient has degree at most $jw$
and norm at most
$\binom hj2^{j\lambda}\le2^{j(\lambda+\ell(h))}$.
For the norm estimate, apply the diagonal similarity specifying
the weighted column norm; the column-norm determinant estimate
bounds each principal minor. In particular
\[
 \deg a_j\le jt,\qquad\log_2\|a_j\|_1\le ja_l.
\]
The constant term cancels in $E=c-h_r(T)$, leaving exactly
$r$ terms, each of weight at most $rt$ and operator norm
at most $2^{ra_l}$. Thus
\[
 \operatorname{wt}(E)\le rt,\qquad
 \log_2\|M_E\|\le ra_l+\ell(b).
\]
The discriminant is weighted homogeneous of weight $d_i(d_i-1)$
in its roots. Its Sylvester determinant has at most $(2d_i-1)!$
terms, each with at most $d_i$ derivative coefficients, bounded
in addition by $d_i$. For a scalar native polynomial, the weighted
last column of its companion matrix gives
$\|[Z^j]f_i\|_1\le2^{(d_i-j)\gamma_i}$ and
$\deg[Z^j]f_i\le(d_i-j)w_i$ for $j<d_i$.
The Sylvester expansion with these coefficient bounds gives degree
at most $u_i$ and logarithmic norm at most $v_i$.

Products and sums now yield
\begin{align*}
 \deg D&\le2rt+2u,&\log_2\|D\|_1&\le2ra_l+2v,\\
 \deg K&\le S_rR_d+2rt+2u,&
 \log_2\|K\|_1&\le S_r(1+R_l)+2ra_l+2v+p+1,\\
 \deg C&\le(S_r+1)R_d+2rt+2u,&
 \log_2\|C\|_1&\le S_r(1+R_l)+2ra_l+2v+p+4+R_l.
\end{align*}
These bounds for $C$ dominate the corresponding bounds for $D$.
Applying them to~\eqref{eq:beta} gives
\[
 \operatorname{wt}(\beta_r)\le W_1r+W_0,\qquad
 \log_2\|M_{\beta_r}\|\le L_1r+L_0.
\]
The principal-minor estimate applied to $H_r$ gives
\[
 \deg H_r\le n(W_1r+W_0+1),\qquad
 \log_2\|H_r\|_1\le n(L_1r+L_0+\ell(n)+1).
\]
Both estimates dominate those for all tail coefficients $a_j$
in $Q_r$: one may directly compare with $bt$ and $ba_l$,
using $n\ge b$, $r\ge1$, and the displayed positive majorants.
For $r<b$, square and sum at most $b+1$ terms; for $r=b$,
keep the bound for $H_b$ without squaring. Sum the bounds over
the factors and use
\[
 2\sum_{r=1}^{b-1}r+b=b^2,\qquad
 2\sum_{r=1}^{b-1}1+1=2b-1.
\]
This proves~\eqref{eq:degree} and~\eqref{eq:height}.
\end{proof}

\section{The arithmetic sources}

The sources in~\cite{Source27,Source29} retain respectively
26 and 28 rational base coordinates, including the distinguished
input. Each has nine independent scalar monic quartics for the
point charts, followed by three or four monic quadratic relations.
Every quadratic radicand belongs to the point algebra $B$.
Each pole generator occurs with exponent exactly one in the complete
guard, whose remaining factor $g$ belongs to $B$. There are no
residual equations or additional indexed guards. These assertions
are checked on every coefficient dependency in the supplied circuits.

For completeness, the separability of the point charts holds at
\emph{every} rational base point. For rational $h_0,h_1$, put
\[
 A_0=h_0^4+12h_0^2h_1^2+4h_1^4-11,\quad
 B_0=4h_0h_1(h_0^2+2h_1^2)+6,
\]
\[
 U_0=6A_0^3-9A_0B_0^2-4B_0^3,\quad
 V_0=-4A_0^3+6A_0B_0^2+3B_0^3.
\]
The point polynomial is
$f(Z)=Z^4-4U_0Z^2+8V_0^2$, with discriminant
\[
 2^{15}V_0^2(2A_0^2-B_0^2)^4(A_0^2-2B_0^2)^2.
\]
The pair $(A_0,B_0)$ cannot vanish: otherwise
$(h_0+h_1\sqrt2)^4=(3-\sqrt2)^2$, whose norm would make
49 a rational fourth power. Neither quadratic norm factor
can vanish at a nonzero rational pair. Finally $V_0=0$ would
produce a rational root of $4X^3-6X-3$, which is Eisenstein
at 3. Thus all hypotheses of Theorem~\ref{thm:main} hold on
every rational base point; no branchwise guard is replaced by
the stronger assertion that its algebra norm is nonzero.

\begin{center}
\begin{tabular}{lrr}
\toprule
Quantity & 27 variables & 29 variables\\
\midrule
$m$ & 12 & 13\\
$b$ & $2^{18}$ & $2^{18}$\\
$n$ & $2^{21}$ & $2^{22}$\\
$w_G$ & $4\,423\,559$ & $6\,789\,625$\\
$w_g$ & $803\,872$ & $1\,234\,707$\\
$\eta$ at exponent $137731$ & $6\,998\,669\,574$ & $7\,005\,250\,004$\\
$\eta_g$ at exponent $137731$ & $1\,106\,350\,691$ & $1\,107\,547\,330$\\
Number of factors & $262\,144$ & $262\,144$\\
\bottomrule
\end{tabular}
\end{center}
The nine point generator weights are 6 and their logarithmic
operator bounds are 14. The remaining weights are
\[
 (1568138,1935231,116318)\quad\text{for 27 variables},
\]
\[
 (1568138,1935231,1935231,116318)\quad\text{for 29 variables}.
\]
Every logarithmic generator bound and integer recurrence is
recorded in the accompanying coefficient specifications.

\begin{corollary}\label{cor:27}
Under the arithmetic premises of the 27-variable source,
the construction gives a computable family in 27 total variables
such that every sufficiently large odd fixed coefficient exponent
defines $\Z$ over $\Q$, and
\[
 \deg P\le111295944630211471361441792<2^{87}.
\]
The sufficient coefficient threshold is ineffective.
\end{corollary}
\begin{proof}
Theorem~\ref{thm:main} preserves the source projection exactly.
The native degree weights are independent of the fixed coefficient
exponent. Substitute them into~\eqref{eq:degree} and apply the
source arithmetic theorem, with its unchanged threshold.
\end{proof}

\begin{corollary}\label{cor:29}
Under the arithmetic premises of the explicit 29-variable source,
the construction at fixed exponent $137731$ defines $\Z$ over
$\Q$ and satisfies
\[
 \deg P\le394212872202031272474705920<2^{89},
\]
\[
 \log_2 H(P)\le558578343639412157650224021485<2^{99}.
\]
\end{corollary}
\begin{proof}
Apply Theorem~\ref{thm:main} to the effective source and substitute
its integer data into~\eqref{eq:degree} and~\eqref{eq:height}.
No arithmetic parameter is enlarged.
\end{proof}

At the same numerical exponent the 27-variable member satisfies
$\log_2 H(P)\le217401420128398649581869203437<2^{98}$.
That exponent is an algebraic reproduction fixture for this family;
it is not certified to exceed the ineffective defining threshold.
Finite algebraic verification does not establish a source arithmetic
premise or make that threshold effective.

\section{Exact specification and verification}

The program \nolinkurl{build_active_specification.py} reads the
pinned integer coefficient circuits, verifies the guard and dependency
assertions, recomputes the weighted companion bounds, and writes
the formulas~\eqref{eq:rootbound}--\eqref{eq:product} with their
explicit finite ranges. The separate checker reconstructs the
coefficient dependencies and integer majorants without importing
the builder. It sums the degree and coefficient estimates over
all $262144$ cases literally, and checks the source and output
hashes and the precise effective or ineffective interpretation.

The independent program \nolinkurl{check_active_degenerations.py}
checks small exact algebras with repeated nonzero $T$-values,
inactive nilpotents, nonreal roots, no rational characters, and
an empty guarded part. It computes each characteristic polynomial
both by determinant expansion and by Newton traces, checks the
idempotent identity, and verifies squarefreeness by an exact
polynomial gcd. These examples test degeneracies; the proof of the
general projection theorem is given above.

The rank-sized matrices and full monomial list are not materialized.
An expression with determinants, coefficient extraction, finite
integer sums and products specifies every coefficient, without
introducing any existential coordinate for an internal index or
matrix entry. All verification programs use the Python standard library.

\begin{thebibliography}{9}
\bibitem{Source27}
\emph{A defining polynomial in twenty-seven variables:
an existence theorem with a uniform degree bound},
companion paper and guarded source in the accompanying repository,
\nolinkurl{papers/few_variables27}.
\bibitem{Source29}
\emph{An explicit twenty-nine-variable polynomial definition of the integers},
companion paper and guarded source in the accompanying repository,
\nolinkurl{papers/few_variables}.
\bibitem{Weighted}
\emph{Weighted bounds for eliminating finite algebraic choices},
the companion-column and principal-minor bounds,
companion paper in \nolinkurl{papers/finite_algebra}.
\end{thebibliography}
\end{document}
